% -*- mode: latex -*- %

\section{Power tuning}
\label{sec:setting-lambda}

The last component of \ppipp is to adaptively tune the weighting parameter
$\lambda$ and thus, in addition to computational benefits, to achieve
statistical efficiency improvements over ``standard'' prediction-powered
inference and classical inference.
Towards this goal, we show how to optimally tune the weighting parameter
$\lambda$.
%
Along the way we make a few theoretical remarks
on the consequent statistical efficiency (see also the
note~\cite{AngelopoulosDuZr23w}).
%
While we cannot expect \ppipp to guarantee
full statistical efficiency for any distribution---this would require
modeling assumptions that
\ppipp explicitly avoids by providing a wrapper around black-box
prediction models $f$---we demonstrate
that it always improves over classical inference and in some cases is
fully statistically efficient.

\subsection{The optimal weighting parameter $\lambda$}

Our first step is to derive the value $\lambda\opt$ of $\lambda$ minimizing
the asymptotic variance $\Tr (\Sigma^\lambda)$ of the prediction-powered
estimate, where we recall the asymptotic covariance
\begin{equation*}
  \Sigma^\lambda
  = H_{\theta\opt}^{-1}
  \left(r \lambda^2 \cov(\nabla \loss_{\theta\opt}^f)
  + \cov(\nabla \loss_{\theta\opt} - \lambda \nabla \loss_{\theta\opt}^f)
  \right) H_{\theta\opt}^{-1}
\end{equation*}
from definition~\eqref{eqn:sigma-lambda}. One could straightforwardly choose
alternative scalarizations of $\Sigma^\lambda$, such as the maximum diagonal
entry $\max_j \Sigma_{jj}^\lambda$, to yield different minimizing
$\lambda$, but the principles remain the same.
(See Appendix~\ref{sec:proof-optimal-lambda} for the quick proof.)

%% We derive $\lambda\opt$ based on the expression for $\Sigma^\lambda$ from Theorem \ref{theorem:normality}. Since $\thetaPP_\lambda$ could be multidimensional, we define $\lambda\opt$ as the minimizer of the cumulative variance:
%% \begin{equation}
%% \label{eq:optimization_problem_lambda}  
%% \lambda\opt = \argmin_{\lambda} \sum_{i=1}^d \Sigma_{ii}^\lambda.
%% \end{equation}
%% Analytically solving for $\lambda\opt$ in the problem \eqref{eq:optimization_problem_lambda} gives the optimal tuning. Below, $\Tr$ denotes the matrix trace.

\begin{proposition}
  \label{prop:optimal-lambda}
  The prediction-powered estimate $\thetaPP_\lambda$ has minimal asymptotic
  variance for
  \begin{equation*}
    \lambda\opt
    \defeq \argmin_\lambda \Tr(\Sigma^\lambda)
    = \frac{\Tr(H_{\theta\opt}^{-1}(
      \Cov(\nabla \loss_{\theta\opt}, \nabla \loss_{\theta\opt}^f)
      + \Cov(\nabla \loss_{\theta\opt}^f, \nabla \loss_{\theta\opt}))
      H_{\theta\opt}^{-1})}{
      2(1+r) \Tr(H_{\theta\opt}^{-1}\Cov(\nabla \loss_{\theta\opt}^f) H_{\theta\opt}^{-1} )}.
  \end{equation*}
\end{proposition}

A few remarks are in order here; we defer discussion of empirical implications
of this tuning to the next section and focus on analysis
here.
Substituting $\lambda\opt$ into the
expression~\eqref{eqn:sigma-lambda} for $\Sigma^\lambda$ yields optimal
asymptotic variance
\begin{equation*}
  \Tr(\Sigma^{\lambda\opt})
  = \Tr(H_{\theta\opt}^{-1}\Cov\left(\nabla \loss_{\theta\opt}\right)
  H_{\theta\opt}^{-1})
  - \frac{\Tr(H_{\theta\opt}^{-1}(\Cov(\nabla \loss_{\theta\opt},
    \nabla \loss_{\theta\opt}^f ) + \Cov(\nabla \loss_{\theta\opt}^f, \nabla \loss_{\theta\opt} ) )H_{\theta\opt}^{-1})^2}{4(1+r) \Tr(H_{\theta\opt}^{-1}\Cov(\nabla \loss_{\theta\opt}^f) H_{\theta\opt}^{-1} ) }.
\end{equation*}
The first term is the cumulative variance of the classical estimate, while
the second term is non-negative. Therefore, we see that with the optimal
choice of $\lambda$, \ppipp per
Algorithm~\ref{alg:general_ppi} (asymptotically) dominates classical
inference. Furthermore, the improvement over classical inference grows with
the correlation of the true gradients and those with predicted labels. By the Cauchy-Schwarz inequality and the law of total variance,
we have
\begin{align*}
  \Tr(\Sigma^{\lambda\opt})
  & \ge \Tr\left(H_{\theta\opt}^{-1} \Cov(\nabla \loss_{\theta\opt})
  H_{\theta\opt}^{-1}\right)
  - \frac{1}{1 + r}
  \Tr\left(H_{\theta\opt}^{-1} \Cov(\E[\nabla \loss_{\theta\opt} \mid X])
  H_{\theta\opt}^{-1}\right) \\
  & =
  \frac{r}{1 + r}
  \Tr\left(H_{\theta\opt}^{-1} \Cov(\nabla \loss_{\theta\opt}) H_{\theta\opt}^{-1}
  \right)
  + \frac{1}{1 + r} \Tr\left(H_{\theta\opt}^{-1}\E[\Cov(\nabla\loss_{\theta\opt} \mid X)]
  H_{\theta\opt}^{-1}\right),
\end{align*}
with equality in the first step if and only if
$\nabla \loss_{\theta\opt}(x, f(x)) = c\E[\nabla \loss_{\theta\opt}(x, Y)
  \mid X = x]$ for some fixed $c \ge 0$. The last bound is
the minimal asymptotic variance of \emph{any} estimator given
both a labeled and unlabeled sample (see~\cite{AngelopoulosDuZr23w}).
Essentially, \ppipp is always more efficient than classical estimators,
and it \emph{can} achieve optimal statistical efficiency if the black-box prediction function $f$ is sufficiently strong.

\subsection{A plug-in estimator and tuning procedure}

In finite samples, $\lambda\opt$ admits the natural
plug-in estimate
\begin{equation}
\label{eq:plugin_lambda}
\hat \lambda = \frac{\Tr(\widehat H_{\thetaPP}^{-1}( \widehat \Cov_n(\nabla \loss_{\thetaPP}, \nabla \loss_{\thetaPP}^f ) + \widehat \Cov_n(\nabla \loss_{\thetaPP}^f, \nabla \loss_{\thetaPP} ) ) H_{\thetaPP}^{-1} )}{2(1+\frac n N) \Tr( \widehat H_{\thetaPP}^{-1}\widehat\Cov_{N+n}(\nabla \loss_{\thetaPP}^f) \widehat H_{\thetaPP}^{-1} ) },
\end{equation}
where $\thetaPP = \thetaPP_\lambda$ is any estimate obtained by taking a
fixed $\lambda \in [0, 1]$, as each of these is $\sqrt{n}$-consistent
(Theorem~\ref{theorem:normality}), and $\what{H}_\theta = \nabla^2
L_n(\theta)$ is the empirical Hessian. The estimate $\hat\lambda$ is
consistent, satisfying
$\hat\lambda \cp \lambda\opt$, as (when the losses are smooth
enough as in Definition~\ref{definition:smooth-enough}) each individual term
consistently estimates its population counterpart.
Theorem~\ref{theorem:normality} thus implies the following corollary.
\begin{corollary}
  \label{cor:clt_w_lambdahat}
  Let $\hat\lambda$ be the estimate \eqref{eq:plugin_lambda} and
  the conditions of Theorem~\ref{theorem:normality} hold. Then
\[\sqrt{n}(\thetaPP_{\hat\lambda}- \theta\opt)\cd \normal(0,\Sigma^{\lambda\opt}).\]
\end{corollary}

The astute reader may notice that both $\lambda\opt$ and $\hat{\lambda}$ may
fail to lie on the region $[0, 1]$, which is the only region in which we can
guarantee convexity---and thus tractability of computing $\thetaPP_{\hat\lambda}$---even for GLMs. (This may occur, for example, if the
predictions $f(X_i)$ are anti-correlated with the true outcomes $Y_i$;
see Example~\ref{example:mean-tuning}.)  In
this case, we can take one of two approaches: the first, naively, is to
simply clip $\hat{\lambda}$ to lie in $[0, 1]$ in Algorithm~\ref{alg:glms}
or \ref{alg:general_ppi}. The second is to use classical one-step
estimators~\cite[e.g.][Chapter 5.7]{VanDerVaart98}. In this case, we
recognize that Theorem~\ref{theorem:normality} implies that for any fixed
$\lambda$, $\thetaPP_\lambda$ is $\sqrt{n}$-consistent, and take a single
Newton step on $\LPP_{\hat\lambda}$,
\begin{equation*}
  \thetaPP_{\textup{os}}
  \defeq \thetaPP_\lambda - \nabla^2 \LPP_{\hat{\lambda}}(
  \thetaPP_\lambda)^{-1} \nabla \LPP_{\hat\lambda}(\thetaPP_\lambda).
\end{equation*}
We call $\thetaPP_{\textup{os}}$ and the resulting inferences ``one-step
PPI.''
%
The difference between the one-step estimator and $\thetaPP_{\hat\lambda}$,
where we clip $\hat\lambda$ to be in $[0, 1]$, is usually empirically
immaterial; we explore this in Appendix~\ref{app:one-step-ppi}. We have the
following corollary~\cite[Thm.~5.45]{VanDerVaart98}.
\begin{corollary}
\label{cor:one-step}
  Let $\hat{\lambda}$ be the estimate~\eqref{eq:plugin_lambda}
  and the conditions of Theorem~\ref{theorem:normality} hold. Then
  \begin{equation*}
    \sqrt{n} (\thetaPP_{\textup{os}} - \theta\opt)
    \cd \normal(0, \Sigma^{\lambda\opt}).
  \end{equation*}
\end{corollary}
Using the plug-in estimate \eqref{eq:plugin_lambda} thus yields a
prediction-powered point estimate with minimal asymptotic variance, which in
turn implies smaller confidence intervals than both PPI according to  problem~\eqref{eq:ppi} and classical inference.

%% As a side note, we generally clip $\hat\lambda$ to be in $[0,1]$ to ensure that the final problem is convex.
%% This is usually irrelevant, since $\hat\lambda < 0$ only when the predictions are far off from the true labels (in the sense that the corresponding gradients are anticorrelated).
%% Furthermore, the clipping is not strictly necessary; we propose a version of the procedure in Section~\ref{app:one-step-ppi} that avoids this issue by taking only a single Newton step.

\begin{example}[Optimal tuning for mean estimation]
  \label{example:mean-tuning}
  The mean-estimation objective with \ppipp is always
  convex for any choice of $\lambda \in \R$, thus obviating the need
  for clipping or one-step corrections.
  Let $\theta\opt = \E[Y] = \argmin_\theta \E[\loss_\theta(Y)]$
  for the loss $\loss_\theta(y) = (y-\theta)^2$.
  Then for any $\lambda \in \R$, we have
  $\thetaPP_\lambda = \frac{1}{n} \sum_{i = 1}^n Y_i
  + \lambda(\frac{1}{N} \sum_{i = 1}^N f(\wt{X}_i) - \frac{1}{n}
  \sum_{i = 1}^n f(X_i))$.
  Substituting into the
  expression for $\lambda\opt$ from Proposition \ref{prop:optimal-lambda}
  gives
  \begin{equation*}
    \lambda\opt = \frac{\Cov(Y,f(X))}{(1+r)\Var(f(X))}.
  \end{equation*}
  We use the plug-in estimator
  \begin{equation*}
    \hat\lambda = \frac{\widehat\Cov_n(Y_i,f(X_i))}{(1+\frac n N )\widehat \Var_{N+n}(f(X_i))}
  \end{equation*}
   of $\lambda\opt$ in our experiments.
\end{example}

%% \subsection{On statistical efficiency}

%% Finally, we turn to a brief discussion of asymptotic efficiency, where we
%% must address a few subtleties. While our estimators are always guaranteed to
%% be more statistically efficient than both classical M-estimators (without
%% unlabeled data) and prediction-powered inference (PPI)---as these correspond
%% to the choices $\lambda = 0$ and $\lambda = 0$, respectively---we do not
%% achieve the full asymptotic statistical efficiency bounds allowed by more
%% sophisticated semiparametric inference
%% techniques~\cite[e.g.][Ch.~25]{VanDerVaart98}. Of course, the point of \ppipp
%% is to provide a lightweight approach that automatically leverages black-box
%% prediction models without sophisticated additional modeling or assumptions,
%% so that we would not expect to achieve full semiparametric
%% efficiency. Nonetheless, we present a few results highlighting what is
%% possible and a few restricted efficiency guarantees.


